% CLEF 2026 Working Note — BirdCLEF+
% CEUR-WS single-column format (standard for CLEF working notes)
\documentclass[a4paper]{article}

\usepackage[utf8]{inputenc}
\usepackage[T1]{fontenc}
\usepackage{amsmath,amssymb}
\usepackage{booktabs}
\usepackage{hyperref}
\usepackage{algorithm}
\usepackage{algorithmic}
\usepackage{multirow}
\usepackage{graphicx}
\usepackage{xcolor}
\usepackage{caption}
\usepackage{subcaption}
\usepackage[margin=2.5cm]{geometry}
\usepackage{enumitem}
\usepackage{natbib}
\bibliographystyle{plainnat}

% ── Custom Commands ──────────────────────────────────────────────────────────
\newcommand{\method}{Noisy Classmate}
\newcommand{\ns}{Noisy Student}
\newcommand{\nc}{Noisy Classmate}
\newcommand{\bce}{\text{BCE}}
\newcommand{\kld}{\text{KLD}}

\title{\textbf{Noisy Classmates: Multi-Architecture Co-Evolutionary\\Self-Training for Bioacoustic Species Recognition}}

\author{
  Tom\\
  Realtek Semiconductor Corp.\\
  Hsinchu, Taiwan\\
  \texttt{tom@realtek.com}
}

\date{BirdCLEF+ 2026 Working Note — CLEF 2026}

\begin{document}
\maketitle

% ═══════════════════════════════════════════════════════════════════════════════
\begin{abstract}
We present \emph{Noisy Classmates} (NC), a multi-architecture co-evolutionary self-training framework for bioacoustic species recognition.
Unlike standard Noisy Student (NS) self-training, where a single model iteratively re-labels and retrains on its own predictions,
NC maintains multiple heterogeneous backbone chains---EfficientNet-B0 and Pyramid Vision Transformer v2---that \emph{teach each other} through shared pseudo labels.
Our framework introduces five progressive innovations:
(1)~ensemble pseudo-label fusion across architectures,
(2)~per-sample confidence-aware blending via inverse-entropy weighting,
(3)~disagreement mining that upweights training samples where models disagree,
(4)~soft-label distillation preserving inter-class probability structure, and
(5)~bidirectional iterative co-evolution where knowledge flows in both directions between chains.
Starting from a Perch v2 foundation model teacher, our pipeline trains 12 rounds of NS for EfficientNet-B0 and 8 rounds for PVT, then transitions to NC for subsequent rounds.
On the BirdCLEF+ 2026 Kaggle leaderboard, our best submission achieves a public score of \textbf{0.942} using only self-trained models (no external competitor weights), representing a +0.011 improvement over the NS-only baseline (0.931).
We provide extensive ablation studies demonstrating that \emph{ensemble diversity}---not individual model accuracy---is the primary driver of leaderboard performance, and that prediction correlation between ensemble members is a more reliable selection criterion than validation AUC.
\end{abstract}

\textbf{Keywords:} bioacoustics, sound event detection, self-training, pseudo labeling, ensemble diversity, noisy student, co-training, knowledge distillation

% ═══════════════════════════════════════════════════════════════════════════════
\section{Introduction}
\label{sec:intro}

Passive acoustic monitoring (PAM) has emerged as a critical tool for biodiversity assessment, enabling continuous, non-invasive surveying of animal populations across large spatial and temporal scales.
The BirdCLEF+ 2026 challenge~\citep{birdclef2026} requires multi-label classification of 234 species---spanning Aves, Amphibia, Insecta, Mammalia, and Reptilia---from 60-second soundscape recordings, evaluated by macro-averaged ROC-AUC across 12 non-overlapping 5-second windows per file.

\paragraph{The Domain Gap.}
The fundamental challenge is the severe distribution mismatch between available training data:
\emph{point recordings} from eBird/xeno-canto (short, single-species, studio-quality) versus
\emph{soundscape recordings} (long, multi-species, environmental noise, varied equipment).
A model trained solely on point recordings generalizes poorly to soundscapes---it has never encountered multi-species overlap, reverberation, or ambient noise.

\paragraph{Self-Training as Domain Adaptation.}
Noisy Student (NS) self-training~\citep{xie2020self} addresses this gap by using a teacher model to pseudo-label unlabeled soundscapes, then training a student on the combination of labeled and pseudo-labeled data with noise injection.
However, standard NS suffers from \emph{confirmation bias}: each round's pseudo labels are generated by the previous round's model, causing systematic errors to be self-reinforced rather than corrected.

\paragraph{Our Contribution.}
We propose \emph{Noisy Classmates} (NC), which replaces the self-referential NS loop with a \emph{peer-learning} paradigm.
Multiple architecturally heterogeneous models---which we call ``classmates''---collaboratively generate pseudo labels for each other.
Because different architectures have different inductive biases (CNNs excel at local spectral texture; Vision Transformers capture global temporal context), their prediction errors are partially independent.
When classmates disagree on a sample, at least one is likely correct---making disagreement the \emph{source of new knowledge} that neither model could discover alone.

Our key contributions are:
\begin{enumerate}[nosep]
    \item A 5-phase progressive framework that extends NS with cross-architecture pseudo-label fusion, confidence-aware blending, disagreement mining, soft-label distillation, and bidirectional co-evolution.
    \item Empirical demonstration that \emph{prediction correlation}, not validation AUC, is the correct criterion for ensemble member selection---overturning the common assumption that stronger individual models yield better ensembles.
    \item A complete, reproducible pipeline achieving LB 0.942 using only self-trained models.
\end{enumerate}

% ═══════════════════════════════════════════════════════════════════════════════
\section{Related Work}
\label{sec:related}

\paragraph{Self-Training and Pseudo Labeling.}
Self-training~\citep{yarowsky1995unsupervised} iteratively labels unlabeled data using model predictions.
Noisy Student~\citep{xie2020self} improved upon this by injecting noise (data augmentation, dropout, stochastic depth) during student training, achieving state-of-the-art on ImageNet.
The BirdCLEF 2025 first-place solution~\citep{babych2025} applied multi-round NS to bioacoustic SED, introducing SumixFreq augmentation and power-transformed pseudo labels.

\paragraph{Co-Training and Mutual Learning.}
Co-training~\citep{blum1998combining} trains two models on different ``views'' of the data, with each model's high-confidence predictions serving as labels for the other.
Deep Mutual Learning~\citep{zhang2018deep} trains peer networks simultaneously, with each minimizing KL divergence to the other's output distribution.
Born-Again Networks~\citep{furlanello2018born} iteratively distill knowledge across generations.
Our NC framework combines elements of all three: multi-view (heterogeneous architectures), mutual pseudo-labeling, and multi-generational evolution.

\paragraph{Ensemble Diversity.}
Krogh and Vedelsby~\citep{krogh1994neural} proved that ensemble error decomposes as average individual error minus diversity (the ``ambiguity decomposition'').
This implies that maximizing diversity---even at the cost of individual accuracy---can reduce ensemble error.
Our empirical findings strongly confirm this: ensemble configurations with lower prediction correlation consistently achieve higher leaderboard scores, even when comprising individually weaker models.

\paragraph{Sound Event Detection.}
Modern SED systems typically use mel-spectrogram front-ends with convolutional or transformer backbones~\citep{kong2020panns}.
Attention-based pooling mechanisms~\citep{mcfee2018adaptive} and frequency-aware pooling~\citep{nam2021frequency} improve temporal and spectral resolution.

% ═══════════════════════════════════════════════════════════════════════════════
\section{Method}
\label{sec:method}

% ─── 3.1 Overview ────────────────────────────────────────────────────────────
\subsection{Pipeline Overview}

Our pipeline consists of three phases, executed sequentially (Figure~\ref{fig:pipeline}):

\begin{enumerate}[nosep]
    \item \textbf{Teacher Extraction}: A Perch v2~\citep{perch2024} foundation model with a fine-tuned linear head generates initial pseudo labels on all 10,658 unlabeled soundscapes.
    \item \textbf{Noisy Student (NS)}: Two independent backbone chains---EfficientNet-B0~\citep{tan2019efficientnet} (12 rounds) and PVT v2 B0~\citep{wang2022pvt} (8 rounds)---each undergo iterative self-training with noise injection.
    \item \textbf{Noisy Classmate (NC)}: The two chains transition from self-training to peer-training, generating shared pseudo labels through a 5-phase innovation stack.
\end{enumerate}

\begin{figure}[t]
\centering
\small
\begin{verbatim}
Phase A: Noisy Student
  Perch → R0 pseudo → B0 R1→R12 (self-training)
                     → PVT R1→R8 (self-training)

Phase B: Noisy Classmate
  B0 R11 + PVT R8 → blend → PVT R9 (Phase 2+3)
  B0 R11 + PVT R9 → blend → PVT R10 (Phase 2+3+4)
  PVT R10 + B0 R11 → blend → B0 R12 (bidirectional!)
  ... (alternating co-evolution)
\end{verbatim}
\caption{Pipeline overview. NS trains independent chains; NC enables cross-chain knowledge flow.}
\label{fig:pipeline}
\end{figure}

% ─── 3.2 SED Architecture ───────────────────────────────────────────────────
\subsection{SED Model Architecture}

Both backbones share the same SED head architecture:

\begin{equation}
\text{audio} \xrightarrow{\text{MelSpec}} \mathbf{X} \in \mathbb{R}^{3 \times 224 \times T}
\xrightarrow{\text{Backbone}} \mathbf{F}
\xrightarrow{\text{GEMFreqPool}} \mathbf{h}
\xrightarrow{\text{AttHead}} \hat{\mathbf{y}} \in [0,1]^{234}
\end{equation}

\paragraph{Mel Spectrogram.}
20-second audio clips are absmax-normalized and converted to mel spectrograms with $n_\text{mels}=224$, $n_\text{fft}=2048$, $\text{hop}=512$, $f_\text{min}=0$, $f_\text{max}=16\text{kHz}$, Slaney normalization, and HTK mel scale.
The spectrogram is converted to dB scale (top\_db=80), min-max normalized, and expanded to 3 channels.

\paragraph{Backbones.}
\textbf{EfficientNet-B0} (5.3M params): CNN backbone with channel attention, pretrained on ImageNet-JFT. Strong at capturing local spectral texture patterns.
\textbf{PVT v2 B0} (3.4M params): Pyramid Vision Transformer with spatial-reduction attention. Captures global temporal dependencies that CNNs miss.

\paragraph{GEMFreqPool.}
Generalized Mean pooling along the frequency axis with learnable exponent $p$ (initialized to 3.0):
$\text{GEM}(\mathbf{x}) = \left(\frac{1}{F}\sum_{f=1}^{F} x_f^p\right)^{1/p}$.
This adaptively interpolates between average ($p=1$) and max ($p \to \infty$) pooling.

\paragraph{Attention SED Head.}
Following~\citet{nam2021frequency}, we use a dual-branch attention head:
attention weights $\alpha = \sigma(W_\text{att} * \mathbf{h})$ and class logits $\mathbf{c} = W_\text{cls} * \mathbf{h}$,
with output $\hat{y}_k = \sigma\left(\sum_t \alpha_{k,t} \cdot c_{k,t}\right)$.

% ─── 3.3 Noisy Student ──────────────────────────────────────────────────────
\subsection{Noisy Student Training}
\label{sec:ns}

\paragraph{Training Data.}
Each round combines labeled point recordings ($\sim$35K clips) with pseudo-labeled soundscapes ($\sim$10K files $\times$ 12 windows).
Pseudo labels are generated by 5-fold ensemble inference on all soundscapes, followed by a Residual Corrector and power-transform sharpening.

\paragraph{Augmentation Pipeline.}
We apply four augmentation stages sequentially:
\begin{enumerate}[nosep]
    \item \textbf{Wave-level MixUp} ($\lambda=0.5$): Each labeled clip is mixed with a pseudo-labeled soundscape clip. Labels are unioned via element-wise max.
    \item \textbf{SpecAugment}: Frequency masking (width 24) and time masking (width 64).
    \item \textbf{SumixFreq}~\citep{babych2025}: For each frequency bin, randomly select between two mel spectrograms. Applied only to B0 (disabled for PVT per empirical finding).
    \item \textbf{Stochastic Depth}: Drop path rate 0.1 in both backbones.
\end{enumerate}

\paragraph{EMA Weight Inheritance.}
Exponential Moving Average (decay=0.999) is maintained during training.
At each round transition, the new model is initialized from the previous round's EMA checkpoint rather than the best checkpoint, providing smoother knowledge transfer and preventing catastrophic forgetting of early-round learning.

\paragraph{Residual Corrector.}
After each round's inference, a lightweight BiSSM network (d\_model=128, d\_state=16) is trained to predict the residual between SED predictions and Perch teacher predictions.
The corrected output is: $\hat{y}_\text{corr} = \hat{y}_\text{SED} + \alpha \cdot \Delta_\text{BiSSM}$ with $\alpha=0.40$.

\paragraph{Pseudo Label Generation.}
Per-round soft pseudo labels are generated via:
(1) gamma sharpening $p_i \leftarrow p_i^\gamma$ (increasing $\gamma$ per round: 1.0$\to$2.0),
(2) per-class 95th-percentile thresholding, and
(3) for non-Aves species, substitution with Perch-only predictions (SED is unreliable for taxa with $<$13 training clips per species).

\paragraph{Teacher Weight Decay.}
The Perch teacher's contribution to pseudo labels decreases across rounds:
$w_\text{teacher}$: R1=0.50, R2=0.30, R3=0.10, R4+=0.00.
This schedule prevents confirmation bias in early rounds while allowing the student to self-correct in later rounds.

% ─── 3.4 Noisy Classmate ────────────────────────────────────────────────────
\subsection{Noisy Classmate: Cross-Architecture Co-Evolution}
\label{sec:nc}

After NS converges (B0 R12, PVT R8), we transition to \emph{Noisy Classmate} training.
The core innovation is replacing self-referential pseudo labels with \emph{cross-architecture ensemble pseudo labels}, implemented through five progressive phases.

\paragraph{Phase 1: Ensemble Pseudo Labels.}
Instead of each chain using its own predictions for the next round, we blend predictions from both chains:
\begin{equation}
\hat{y}_\text{NC}^{(i)} = w_\text{B0} \cdot \hat{y}_\text{B0}^{(i)} + w_\text{PVT} \cdot \hat{y}_\text{PVT}^{(i)}
\end{equation}
where $w_\text{B0} = w_\text{PVT} = 0.5$. This is analogous to bagging: averaging predictions from models with partially independent errors reduces variance without increasing bias.

\paragraph{Phase 2: Confidence-Aware Blending.}
Phase 1 treats all predictions equally. Phase 2 weights each model's contribution by its per-sample confidence (inverse entropy):
\begin{equation}
w_m^{(i)} = \frac{\text{base}_m / H(\hat{y}_m^{(i)})}{\sum_{m'} \text{base}_{m'} / H(\hat{y}_{m'}^{(i)})}
\end{equation}
where $H(\hat{y}) = -\sum_k [p_k \log p_k + (1-p_k)\log(1-p_k)]$ is the binary entropy over all classes.
This automatically delegates each sample to its ``specialist'' classmate.

\paragraph{Phase 3: Disagreement Mining.}
The most informative training samples are those where classmates \emph{disagree}---analogous to Query-by-Committee active learning~\citep{seung1992query}.
We compute per-sample disagreement:
\begin{equation}
d^{(i)} = \frac{1}{K} \sum_{k=1}^{K} \text{Var}\left(\hat{y}_{1,k}^{(i)}, \hat{y}_{2,k}^{(i)}\right)
\end{equation}
and assign training weights $w_\text{train}^{(i)} = 1 + \alpha \cdot d^{(i)} / \max_j d^{(j)}$ with $\alpha=2.0$ (maximum 3$\times$ upweighting).
Labeled samples receive weight 1.0; only pseudo-labeled samples are disagreement-weighted.

\paragraph{Phase 4: Soft Label Distillation.}
Standard pseudo labels undergo gamma sharpening and thresholding, discarding inter-class similarity information.
Following Hinton et al.~\citep{hinton2015distilling}, we preserve soft probability distributions via a combined loss:
\begin{equation}
\mathcal{L} = (1-\beta) \cdot \mathcal{L}_\text{Focal}(\mathbf{z}, \mathbf{y}_\text{hard})
+ \beta \cdot T^2 \cdot \kld\left(\sigma(\mathbf{z}/T) \,\|\, \mathbf{y}_\text{soft}\right)
\end{equation}
with $\beta=0.3$ and temperature $T=2.0$. The KLD term teaches the student that ``this sound is 60\% species A and 30\% species B,'' rather than just ``this is species A.''

\paragraph{Phase 5: Bidirectional Co-Evolution.}
Phases 1--4 enable one-directional teaching (B0$\to$PVT).
Phase 5 closes the loop: after PVT trains on B0's knowledge and develops its own perspective, its predictions feed \emph{back} into B0's pseudo labels.
This creates an iterative dialogue:
\begin{align}
&\text{Gen } n: \quad \text{PVT}_{r+1} \leftarrow \text{blend}(\text{B0}_R, \text{PVT}_r) \\
&\text{Gen } n: \quad \text{B0}_{R+1} \leftarrow \text{blend}(\text{PVT}_{r+1}, \text{B0}_R)
\end{align}
Each generation, both models learn from the other's unique perspective.
Convergence is guaranteed since architectural inductive biases prevent the models from collapsing to identical predictions.

% ─── 3.5 Inference ───────────────────────────────────────────────────────────
\subsection{Inference and Ensemble}
\label{sec:inference}

\paragraph{ONNX Export.}
Trained checkpoints are exported to ONNX format with optional INT8 dynamic quantization (3.7$\times$ compression with negligible accuracy loss).
Inference runs on CPU via ONNX Runtime in the Kaggle notebook.

\paragraph{Ensemble Selection.}
We select three models for the final ensemble based on empirical rules derived from extensive leaderboard experiments (Section~\ref{sec:diversity}):
\begin{enumerate}[nosep]
    \item Maximize \emph{round diversity}: models from different NS rounds have lower prediction correlation.
    \item Include \emph{architecture diversity}: at least one CNN (B0) and one ViT (PVT).
    \item \emph{Never} select two models from the same NS round.
\end{enumerate}

Our best ensemble: B0 R12 fold0 + PVT R5 fold4 + B0 R6 fold3 (LB 0.942).

\paragraph{VLOM Blend.}
SED ensemble predictions are combined with Perch probe predictions via Variance-weighted Log-Odds Mean:
\begin{equation}
\hat{y}_\text{final} = \sigma\left(w_\text{SED} \cdot \text{logit}(\hat{y}_\text{SED}) + w_\text{Perch} \cdot \text{logit}(\hat{y}_\text{Perch})\right)
\end{equation}
with $w_\text{SED}=0.7$ and $w_\text{Perch}=0.3$.

% ═══════════════════════════════════════════════════════════════════════════════
\section{Experiments}
\label{sec:experiments}

% ─── 4.1 Setup ───────────────────────────────────────────────────────────────
\subsection{Experimental Setup}

\paragraph{Data.}
BirdCLEF+ 2026 provides $\sim$35K labeled point recordings and 10,658 unlabeled 60-second soundscape recordings across 234 species.
66 soundscapes have sparse ground-truth labels (739 labeled windows).
We use 5-fold cross-validation on soundscape files for model selection.

\paragraph{Training.}
All training on a single NVIDIA GPU (24GB).
B0: 30 epochs (R1--R4), 25 epochs (R5+); PVT: 25 epochs (all rounds).
AdamW optimizer, initial LR $10^{-3}$ (B0) or $5 \times 10^{-4}$ (PVT R5+), cosine annealing to $10^{-6}$.
Early stopping with patience 3--4 based on soundscape validation AUC.

% ─── 4.2 NS Results ──────────────────────────────────────────────────────────
\subsection{Noisy Student Results}

Table~\ref{tab:ns_rounds} shows the progression of soundscape validation AUC across NS rounds.

\begin{table}[t]
\centering
\caption{Soundscape validation AUC (5-fold mean) across NS rounds.}
\label{tab:ns_rounds}
\small
\begin{tabular}{lccccc}
\toprule
Round & B0 Mean & Best B0 Fold & PVT Mean & Best PVT Fold \\
\midrule
R1  & 0.9120 & f3: 0.929 & 0.9005 & f1: 0.924 \\
R2  & 0.9256 & f3: 0.940 & 0.9225 & f3: 0.946 \\
R4  & 0.9341 & f3: 0.954 & 0.9409 & f3: 0.957 \\
R6  & 0.9383 & f3: 0.966 & 0.9513 & f2: 0.966 \\
R8  & 0.9383 & f3: 0.960 & 0.9545 & f2: 0.969 \\
R10 & 0.9435 & f0: 0.964 & --- & --- \\
R12 & 0.9468 & f0: 0.965 & --- & --- \\
\bottomrule
\end{tabular}
\end{table}

\paragraph{Observations.}
B0 shows rapid improvement R1--R4 (+0.022) followed by diminishing returns R4--R8 (+0.004) and a distribution shift at R9 (fold0 jumps from 0.90 to 0.96 due to pseudo-label regime change).
PVT improves more consistently, reaching mean 0.9545 at R8.

% ─── 4.3 NC Results ──────────────────────────────────────────────────────────
\subsection{Noisy Classmate Results}

Table~\ref{tab:nc_vs_ns} compares NC R9 against NS R8 for PVT.

\begin{table}[t]
\centering
\caption{NC vs NS: PVT per-fold soundscape AUC comparison.}
\label{tab:nc_vs_ns}
\small
\begin{tabular}{lccccc}
\toprule
 & fold0 & fold1 & fold2 & fold3 & fold4 \\
\midrule
PVT R8 (NS) & 0.9575 & 0.9435 & 0.9693 & 0.9603 & 0.9418 \\
PVT R9 (NC) & \textbf{0.9639} & \textbf{0.9449} & 0.9662 & \textbf{0.9678} & --- \\
$\Delta$     & +0.006 & +0.001 & $-$0.003 & +0.008 & --- \\
\bottomrule
\end{tabular}
\end{table}

NC improves 3 out of 4 completed folds, with the largest gain at fold3 (+0.008). This demonstrates that cross-architecture pseudo labels provide complementary learning signal beyond what self-training can achieve.

% ─── 4.4 Diversity Analysis ──────────────────────────────────────────────────
\subsection{The Diversity Paradox}
\label{sec:diversity}

Our most significant finding is the \emph{Diversity Paradox}: ensemble configurations with lower average validation AUC can achieve \emph{higher} leaderboard scores.

\begin{table}[t]
\centering
\caption{Leaderboard scores vs. ensemble composition. Avg AUC is the mean of individual model soundscape validation AUCs. Corr is prediction correlation between the two B0 models.}
\label{tab:diversity}
\small
\begin{tabular}{llccc}
\toprule
Config & Models & Avg AUC & Corr & LB \\
\midrule
A & B0 R12f0 + PVT R5f4 + B0 R6f3  & 0.960 & 0.978 & \textbf{0.942} \\
B & B0 R12f0 + PVT R5f4 + B0 R8f3  & 0.959 & 0.983 & 0.941 \\
C & B0 R12f0 + PVT R5f4 + B0 R12f2 & 0.958 & 0.987 & 0.940 \\
D & B0 R8f2 + PVT R5f4 + B0 R8f3   & 0.947 & 0.983 & 0.938 \\
E & B0 R9f0 + PVT R8f2 + B0 R6f3   & 0.967 & --- & 0.937 \\
\bottomrule
\end{tabular}
\end{table}

\paragraph{Key Finding 1: Correlation predicts LB better than AUC.}
Configs A, B, C differ only in the third model. Their LB ranking (0.942 $>$ 0.941 $>$ 0.940) perfectly anti-correlates with prediction correlation (0.978 $<$ 0.983 $<$ 0.987), but does \emph{not} correlate with average AUC.

\paragraph{Key Finding 2: Round diversity drives decorrelation.}
Config A uses R6 and R12 (6 rounds apart); B uses R8 and R12 (4 apart); C uses R12 and R12 (same round).
Models from different NS rounds were trained on different pseudo labels, producing different error patterns.

\paragraph{Key Finding 3: The GT Paradox.}
Config E has the highest average AUC (0.967) but the lowest LB (0.937).
All three models have individually strong but \emph{uniformly high} AUCs (0.965--0.969), indicating high correlation.
This confirms the ambiguity decomposition~\citep{krogh1994neural}: ensemble error = avg error $-$ diversity.

\paragraph{Ensemble Selection Algorithm.}
Based on these findings, we propose a diversity-first selection:
\begin{enumerate}[nosep]
    \item Fix the PVT model (R5 fold4, the only effective PVT configuration).
    \item Select B0 model 1 from a late round (R12) for strong individual accuracy.
    \item Select B0 model 2 from an \emph{early} round (R4--R6) to maximize round diversity.
    \item Verify prediction correlation $<$ 0.985 between the two B0 models.
\end{enumerate}

% ─── 4.5 Ablation ────────────────────────────────────────────────────────────
\subsection{Ablation Studies}

\begin{table}[t]
\centering
\caption{Ablation of key design choices on LB score.}
\label{tab:ablation}
\small
\begin{tabular}{lc}
\toprule
Configuration & LB \\
\midrule
\textbf{Full pipeline (best)} & \textbf{0.942} \\
\quad $-$ Round diversity (use R12+R12 instead of R12+R6) & 0.940 \\
\quad $-$ Architecture diversity (B0-only, no PVT) & 0.938 \\
\quad $-$ fold0 (use fold2 for first B0) & 0.938 \\
\quad $-$ VLOM (equal blend sed/perch 0.5/0.5) & 0.937 \\
\quad $-$ All models at R5 (no later rounds) & 0.933 \\
\quad $-$ NS entirely (use Perch-only) & 0.915 \\
\bottomrule
\end{tabular}
\end{table}

\paragraph{Impact Ranking.}
The largest single contributor is multi-round NS training itself (+0.018 over Perch-only).
Within the ensemble, architecture diversity (+0.004) and round diversity (+0.002) provide the next largest gains.

% ═══════════════════════════════════════════════════════════════════════════════
\section{Analysis}
\label{sec:analysis}

\subsection{Why NC Works: Breaking Confirmation Bias}

Standard NS suffers from confirmation bias because the model generates its own training labels.
If B0 R5 systematically misclassifies a rare species, this error propagates to R6--R12.
NC breaks this cycle: PVT's different architecture may correctly identify the species that B0 misses.
When the blend includes PVT's correct prediction, B0's successor learns the correct label.

\subsection{Why Different Architectures Help}

CNNs (EfficientNet-B0) process spectrograms with local receptive fields, excelling at spectral texture patterns (e.g., whistle-type bird calls with clear frequency contours).
Vision Transformers (PVT v2) use global self-attention, capturing temporal dependencies across the full spectrogram (e.g., rhythmic insect calls, frog choruses).
These complementary strengths mean their errors are partially independent---satisfying the theoretical requirements of co-training~\citep{blum1998combining}.

\subsection{Convergence Properties}

NC co-evolution will converge when both models' predictions become identical (diversity $\to$ 0).
However, because architectural inductive biases are fundamentally different, full convergence is impossible---residual diversity always exists.
In practice, we observe diminishing returns after 2--3 generations of co-evolution, suggesting an optimal stopping point.

% ═══════════════════════════════════════════════════════════════════════════════
\section{Conclusion}
\label{sec:conclusion}

We presented \emph{Noisy Classmates}, a framework that extends Noisy Student self-training with cross-architecture peer learning.
Our five-phase innovation stack---ensemble pseudo labels, confidence-aware blending, disagreement mining, soft-label distillation, and bidirectional co-evolution---enables models with different inductive biases to teach each other, creating knowledge that neither could discover alone.

Our extensive leaderboard experiments revealed a counterintuitive finding: \emph{ensemble diversity matters more than individual model accuracy}.
Prediction correlation between ensemble members is a more reliable selection criterion than validation AUC, and round diversity (using models from different self-training rounds) is a practical way to maximize decorrelation.

The complete pipeline achieves LB 0.942 on BirdCLEF+ 2026 using only self-trained models, with no reliance on external competitor weights.
All code, configurations, and reproduction specifications are publicly available.

% ═══════════════════════════════════════════════════════════════════════════════
% References
\begin{thebibliography}{20}

\bibitem[Babych(2025)]{babych2025}
Babych, V. (2025).
1st Place Solution to BirdCLEF 2025.
\textit{Kaggle Discussion}.

\bibitem[BirdCLEF(2026)]{birdclef2026}
BirdCLEF+ 2026.
\textit{Kaggle Competition}.
\url{https://www.kaggle.com/competitions/birdclef-2026}

\bibitem[Blum \& Mitchell(1998)]{blum1998combining}
Blum, A. and Mitchell, T. (1998).
Combining Labeled and Unlabeled Data with Co-Training.
\textit{COLT}.

\bibitem[Furlanello et al.(2018)]{furlanello2018born}
Furlanello, T., Lipton, Z., Tschannen, M., Itti, L., and Anandkumar, A. (2018).
Born Again Neural Networks.
\textit{ICML}.

\bibitem[Hinton et al.(2015)]{hinton2015distilling}
Hinton, G., Vinyals, O., and Dean, J. (2015).
Distilling the Knowledge in a Neural Network.
\textit{NeurIPS Workshop}.

\bibitem[Kong et al.(2020)]{kong2020panns}
Kong, Q., Cao, Y., Iqbal, T., Wang, Y., Wang, W., and Plumbley, M. D. (2020).
PANNs: Large-Scale Pretrained Audio Neural Networks for Audio Pattern Recognition.
\textit{IEEE/ACM TASLP}.

\bibitem[Krogh \& Vedelsby(1994)]{krogh1994neural}
Krogh, A. and Vedelsby, J. (1994).
Neural Network Ensembles, Cross Validation, and Active Learning.
\textit{NeurIPS}.

\bibitem[McFee et al.(2018)]{mcfee2018adaptive}
McFee, B., Salamon, J., and Bello, J. P. (2018).
Adaptive Pooling Operators for Weakly Labeled Sound Event Detection.
\textit{IEEE/ACM TASLP}.

\bibitem[Nam et al.(2021)]{nam2021frequency}
Nam, H., Kim, S. H., and Park, J. (2021).
Frequency Dynamic Convolution: Frequency-Adaptive Pattern Recognition for Sound Event Detection.
\textit{Interspeech}.

\bibitem[Perch(2024)]{perch2024}
Google Research. (2024).
Perch: A Large-Scale Bird Sound Recognition Model.

\bibitem[Seung et al.(1992)]{seung1992query}
Seung, H. S., Opper, M., and Sompolinsky, H. (1992).
Query by Committee.
\textit{COLT}.

\bibitem[Tan \& Le(2019)]{tan2019efficientnet}
Tan, M. and Le, Q. V. (2019).
EfficientNet: Rethinking Model Scaling for Convolutional Neural Networks.
\textit{ICML}.

\bibitem[Wang et al.(2022)]{wang2022pvt}
Wang, W., Xie, E., Li, X., Fan, D., Song, K., Liang, D., Lu, T., Luo, P., and Shao, L. (2022).
PVT v2: Improved Baselines with Pyramid Vision Transformer.
\textit{Computational Visual Media}.

\bibitem[Xie et al.(2020)]{xie2020self}
Xie, Q., Luong, M., Hovy, E., and Le, Q. V. (2020).
Self-Training with Noisy Student Improves ImageNet Classification.
\textit{CVPR}.

\bibitem[Yarowsky(1995)]{yarowsky1995unsupervised}
Yarowsky, D. (1995).
Unsupervised Word Sense Disambiguation Rivaling Supervised Methods.
\textit{ACL}.

\bibitem[Zhang et al.(2018)]{zhang2018deep}
Zhang, Y., Xiang, T., Hospedales, T. M., and Lu, H. (2018).
Deep Mutual Learning.
\textit{CVPR}.

\end{thebibliography}

\end{document}
